\begin{frame}{왜 ``가까워질수록 더 멀어지는가''}
  \begin{itemize}
    \item 1.0 $\to$ 1.495 $\to$ 1.891은 \textbf{자기 가속} 구조.
    \item $\theta$가 목표에 가까워질수록 $\theta$가 커지고,
      $\theta$가 커질수록 목표는 더 멀리 도망간다.
    \item 스텝 크기가 충분히 크면 이 \textbf{양(+) 피드백}이
      \textbf{진동 또는 발산}으로 커진다.
  \end{itemize}
  \vskip0.5em
  \begin{block}{초기 경험의 주된 원인}
    ``DQN을 돌려 보면 가끔은 잘 되고, 가끔은 손실이 폭주한다''는
    바로 이 메커니즘.
  \end{block}
  \begin{itemize}
    \item 타겟 네트워크가 있는 경우: $\theta^{-}$ 동기화라는
      ``벽''이 $C$스텝마다 한 장씩 쳐지므로, 그 사이 목표는
      $\theta$의 현재값과 \textbf{무관하게} 고정.
  \end{itemize}
\end{frame}
